% Pendiente de desarrollo.

\clearpage
\subsection*{Declaración sobre el uso de inteligencia artificial generativa}
\addcontentsline{toc}{subsection}{Declaración sobre el uso de inteligencia artificial generativa}

En coherencia con los principios de transparencia académica y con las recomendaciones emitidas por la UNESCO, la Comisión Europea y un número creciente de universidades españolas sobre el uso responsable de sistemas de inteligencia artificial generativa en trabajos académicos, se declara expresamente la utilización de este tipo de herramientas durante el desarrollo del presente Trabajo Fin de Máster. Esta declaración tiene por objeto explicitar en qué fases del trabajo se han empleado, con qué finalidad, y de qué manera se ha asegurado que su uso no comprometa ni la autoría intelectual, ni el rigor metodológico, ni la validez de los resultados.

La distinción fundamental adoptada en este trabajo diferencia entre los modelos de inteligencia artificial que constituyen el objeto de estudio (y que por tanto forman parte integral de la metodología descrita) y las herramientas de IA generativa utilizadas como apoyo instrumental durante la elaboración del trabajo.

\subsubsection*{Modelos de IA como componentes metodológicos del estudio}

El presente trabajo se apoya, por su propia naturaleza, en varios modelos de inteligencia artificial que constituyen la columna vertebral del \textit{pipeline} analítico y cuya utilización no se considera «uso de IA generativa» en sentido estricto, sino aplicación de técnicas de NLP ampliamente validadas en la literatura científica:

\begin{itemize}
    \item \textbf{all-MiniLM-L6-v2} (Sentence-Transformers): \textit{embeddings} para BERTopic.
    \item \textbf{facebook/bart-large-mnli}: clasificador \textit{zero-shot}.
\end{itemize}

\subsubsection*{Uso de IA generativa como herramienta de apoyo}

Se ha utilizado un asistente conversacional de IA generativa (Claude, de Anthropic) como herramienta de apoyo para tareas instrumentales y de revisión, bajo supervisión humana directa, en:

\begin{itemize}
    \item \textbf{Apoyo en el desarrollo de código Python:} revisión de \textit{scripts}, identificación de \textit{bugs}, generación de primeras versiones de los \textit{notebooks}. El código final ha sido ejecutado, validado y adaptado por la autora.
    \item \textbf{Orientación metodológica:} discusión sobre BERTopic vs LDA, \textit{zero-shot} vs \textit{fine-tuning}, estrategias de validación.
    \item \textbf{Búsqueda y contraste de cifras actualizadas:} FTC 2024 e INCIBE 2025, verificadas accediendo a las fuentes originales.
\end{itemize}

\clearpage
\subsection*{Taxonomía de roles CRediT}
\addcontentsline{toc}{subsection}{Taxonomía de roles CRediT}

\textbf{Reyes Piñón Gómez:} conceptualización, \textit{data curation}, investigación, metodología, redacción del \textit{draft} inicial, supervisión.

\textbf{Alicia Alcántara Troncoso:} conceptualización, \textit{data curation}, investigación, supervisión.
